\documentclass[11pt,a4paper]{article}
\usepackage[T1]{fontenc}
\usepackage[utf8]{inputenc}
\usepackage[margin=22mm]{geometry}
\usepackage{lmodern}
\usepackage{amsmath,amssymb}
\usepackage{enumitem}
\usepackage{xcolor}
\usepackage{microtype}
\usepackage{hyperref}
\definecolor{epflred}{HTML}{E3000B}
\hypersetup{colorlinks=true,urlcolor=epflred,linkcolor=epflred,
  pdftitle={ENG-654 Project 20: Repeatable Closed-Path Motion on iiwa 7}}
\setlength{\parindent}{0pt}
\setlength{\parskip}{6pt}
\setlist[itemize]{leftmargin=*,itemsep=5pt,topsep=4pt}

\begin{document}
{\small\textbf{EPFL \textbar\ ENG-654}\hfill Kinematics-Grounded Motion Planning for Robots}
\vspace{5mm}

{\color{epflred}\Large\bfseries Project 20}

{\LARGE\bfseries Repeatable Closed-Path Motion on iiwa 7\par}

\textbf{Format:} Group of two students; simulation-based project.\\
\textbf{Prerequisites:} Closed-path continuation, redundant inverse kinematics, null-space motion, and graph-based path planning.

\section*{Motivation}
A redundant robot can return its tool to the starting pose while
ending in a different posture. Repeating the same task can then accumulate
joint drift or encounter new constraints. Repeatable motion requires
closing the joint trajectory as well as the workspace loop.

\section*{Task}
\textbf{Overall objectives.} Use classified iiwa IK feasibility charts to seek a redundancy
schedule that closes in both full-pose workspace and joint space, and
compare it with local continuation of the same closed task.

\begin{itemize}
  \item \textbf{Prepare an explicitly closed pose loop.} Use the supplied rectangular path
or a short derived loop and verify both translational and rotational
closure. Add a correctly interpolated closing segment if needed; do not
assume the first and last CSV rows are identical. Validate the model,
tool transform, limits, and one admissible initial configuration.

  \item \textbf{Measure a local baseline.} Track the loop with the same
pseudoinverse/null-space policy on every lap. Continue from the actual
final configuration for up to three laps and record endpoint joint
displacement, pose error, and feasibility. A closed tool pose alone does
not establish joint-space repeatability.

  \item \textbf{Generate and classify the global chart.} Evaluate fixed-$q_3$
IK families over one full loop. Plot path sample versus redundancy angle,
with branch labels and failure masks. Inspect the first and last columns
as IK sets of the same pose, including their configuration correspondence.

  \item \textbf{Enforce joint-space closure.} Use a layered search with the
initial configuration fixed and require the last configuration to match
it within a declared tolerance and consistent permitted angle
representatives. Validate the closing edge. Returning to the same
redundancy angle or branch label alone is insufficient.

  \item \textbf{Compare repeated execution.} Refine the selected schedule and
validate all transitions, including the seam between laps. Compare the
local and closed schedules by endpoint joint error, travel, minimum
limit margin, and full-Jacobian regularity. Show both trajectories on
the feasibility chart and explain any accumulated posture drift.
\end{itemize}

\newpage
\section*{Specifics and scope}
Use one closed pose loop, one shared starting configuration, and
at most three validation laps. Begin with about 150 pose samples and
120 redundancy angles. Search only within the selected sampled families;
if no closing schedule is found, report the search limits rather than
claiming global impossibility. Timing optimization is not required.

For every chart, keep pose validity, joint limits,
full-Jacobian regularity, and reduced-chart regularity as separate diagnostics.
The rectangular $6\times7$ Jacobian has no determinant. A feasible pixel is
not a validated edge: re-solve intermediate prescribed poses and check joint
continuity within limits. If collision checking is unavailable, state that
feasibility is kinematic only; a clearance proxy is not a full collision test.

\section*{Expected outcome}
Understand the distinction between tool-pose closure and repeatable
robot motion. Deliver a classified redundancy chart with endpoint IK
matches, a comparison of local and closure-constrained schedules, and
lap-by-lap evidence of joint closure or drift.

Submit reproducible code, model parameters and test data, the required plots,
a concise 5--6-page report, and a short simulation demonstration. Explain
both successful cases and failures; conclusions must follow the numerical
and geometric evidence rather than an expected outcome.

\section*{Resources}
The KUKA LBR iiwa 7 robot description (URDF) and visualization meshes
(STL) will be provided to students, together with a mathematical continuation
model with unrestricted joints and an analytical fixed-$q_3$ IK model. Reuse
the supplied IK rather than deriving a general 7R solver. Apply the physical
joint limits from the robot description separately when using the unrestricted
model. Check D--H parameters, joint axes, and the fixed tool transform against
the URDF, and validate IK outputs by independent forward kinematics.

A sampled rectangular full-pose path and a reference fixed-$q_3$ IK
table will also be provided. The pose data contain position coordinates and
quaternion columns \texttt{qw,qx,qy,qz}; normalize the quaternions and make
their signs consistent before interpolation. The reference IK table checks
solutions at $q_3=30^\circ$ only; it does not replace evaluation at other angles.
Use its configurations to check pose residuals and branch labels, applying
physical joint limits separately.

Implement the local continuation baseline and a layered search through the
redundant IK map using the supplied model and IK. Check workspace closure,
match endpoint joint configurations, and validate the closing edge and
inter-lap continuity before visualizing the trajectories with the supplied
meshes.

\end{document}
